
%COVID cancelled standardized testing in the UK and lead to record number of  top grades getting issues for two years straight, and overstreching universities capacities.
How would teachers and schools grade students differently compared to a standardized grading system where teachers have no discretion? The COVID-19 pandemic canceled standardized exams in the UK for two consecutive years, and students were required to use their teachers’ predicted grades as substitutes. These teacher-assigned grades determined students’ placement outcomes, as the UK higher education system bases admission decisions on applicants’ standardized exam results (e.g., A-levels).The policy led to a 12-percentage-point increase in the number of top grades awarded in 2020, and an even larger increase in 2021, despite the disruptions in learning for the 2021 cohort. Consequently, top universities became overcrowded as they admitted more students than their capacity, while lower-ranked universities suffered from shortages of students. This episode underscores how a change in the grading system significantly reshaped the grade distribution, with far-reaching consequences for the higher education market.

% Why this is economics? Because its localization of information.
I leverage the pandemic-driven transfer of grading discretion from a centralized examiner to local teachers to estimate institutional differences in grading patterns. To identify such differences, researchers compare outcomes under standardized exams with those under teachers’ discretion, aggregate the discrepancies at the school level, and compare them across institutions\footnote{Other papers detect anomalies by comparing students near the cutoff of a passing grade (\cite{neal2008}, \cite{diamond}). Grading differences across demographic groups are estimated by comparing the size of the grading discontinuities across the passing threshold by demographic groups. Such approaches do not compare across exams, and derives grading anomalies using variation across the same test.}. Several studies document systematic differences between teacher-assigned and centrally examined grades (\cite{10.1162/00335530360698441}, \cite{lavy2009}, \cite{hanna2012}, \cite{dee2019}), showing that teachers often grade differently than external examiners. However, these papers do not compare grading policies across schools with different institutional characteristics.\footnote{\cite{figlio2006} examines different responses from schools to external incentives to increase student performance, finding how private schools mis-reports their students disability status to hide the bad performance of lower performing students. However they do not focus on grades of the students as the schools in their setting used standardized tests to measure a schools performance.}. This gap partly reflects a practical limitation: most schools do not administer comparable standardized tests that would allow researchers to benchmark student ability and directly compare grading policies across institutions.

% \cite{10.1257/app.20200525} find evidence of persistence of grade inflation over time.  However, none of the papers document how grading policies across institutions varies across the entire national secondary school system as well their importance relative to grade inflation due to student's characteristics. Additionally, no papers examine the dynamic persistence in grade inflation practices overtime as well as their granular compositions by smaller subgroups such as subject groups. 

% Why this is economics? Grades determines the labor market outcomes for university graduates. 
I study whether the policy-induced increases in students’ grades improved university placements, especially for those from socially and economically disadvantaged backgrounds. The grades assigned during the pandemic were an unprecedented opportunity for students to improve their placements, since universities were legally obliged to accept all students who met their conditional offers. A large literature (\cite{manski1983college}; \cite{Chetty2023}) examines policies to expand access for disadvantaged students — for example, financial aid (\cite{10.1257/000282803321455287}), affirmative action (\cite{10.1093/qje/qjab027}), and application assistance (\cite{10.1093/qje/qjs017}). However, these papers do not examine how schools and teachers use grades strategically to place students at selective institutions.\footnote{\cite{denning2023} examines how grade inflation contributed to rising average high-school GPAs at selective universities.} This paper tests which schools capitalized on the change in the grading regime to improve their students’ placements and evaluates the implications of these practices for the distribution of students across the higher-education system.

% How schools inflate differently. Static patterns as well as their dynamics.
I test whether schools systematically differed in their tendency to assign inflated grades relative to the pre-pandemic period. I compare final grades between cohorts subject to different exam regimes and assess the extent of grade increases among academically and demographically similar students. A primary identification threat is that cross-school divergences in grades could reflect pre-existing trends rather than the change in grading method. This is less of a concern in my setting because standardized tests in the UK (e.g. \emph{A-levels}) are regulated so that the overall grade distribution reflects the historical relationship between students’ entry-level exams (\emph{GCSE}) and final A-level outcomes. I confirm stationary property of the grade distribution through a myriad of stationarity tests. Additionally, I examine how students with different academic backgrounds (race, gender, economic affluence, and GCSE scores) were graded relative to their counterparts before the pandemic, comparing patterns across both pandemic years to capture changes in schools’ grading behavior. Finally, I test whether schools differed in their success rates at placing students into their top-choice programs, controlling for students’ predicted A-level scores and the selectivity of the target programs.

My approach extends existing methods for detecting changes in grading patterns (\cite{neal2008}; \cite{dee2019}), which estimate grading biases by comparing students in different cohorts who were graded under different methods. I contribute to this literature by exploiting quasi-experimental variation in grades generated by the pandemic shock. I also extend the model to allow grading to vary by students’ academic and demographic characteristics and to differ across years. Finally, I examine how both the magnitude of grade inflation and its targeting across student groups affect university-application success rates, and how the margin of conversion from grades to admissions differs across schools.

I find substantial heterogeneity in how schools loosened their grading policies: many granted students higher grades than they would have received under standardized exams. The extent of this loosening varied by school quality and by school type. Schools of lower quality—measured by past cohort performance on A-levels—exhibited larger increases in grade inflation (measured as the rise in a student’s probability of receiving a top grade), whereas schools with higher baseline rates of top grades showed lower increases. I also document differences by school type: private (independent) institutions exhibited the largest upward shifts. Students in independent schools were about 7 percentage points more likely to receive top grades than students in three-year secondary schools (i.e., sixth-form colleges), which experienced the smallest increases. I find most schools revised their grading policies upward in the second pandemic year, further loosening standards and increasing the share of top grades. Private (independent) schools showed the largest upward revisions. Persistence in grading behavior was stronger for non-quantitative subjects, areas where teachers exercise greater discretionary judgment, whereas quantitative subjects exhibited weaker persistence in inflation. Moreover, schools with relatively high inflation in quantitative subjects in the initial year also tended to show larger inflation in non-quantitative subjects that same year.

Within subject courses, top grades were concentrated among students with stronger academic backgrounds: students with higher prior math and English scores were more likely to receive A/A* grades. Grade inflation among top academic performers at high-quality schools was limited, since these students already had a high probability of earning top grades and thus gains were capped. I find that students from socioeconomically disadvantaged backgrounds received higher grades at low-quality schools, whereas mid- to high-quality schools did not systematically grade students differently by socioeconomic background. I also find a persistent gender gap: conditional on prior performance, female students were more likely than male students to receive top grades—among students with a 50\% baseline probability of an A or A*, males were roughly 6 percentage points less likely than females to receive the highest grades—and this result is robust to subject-level controls. By contrast, racial differences were small: white students were about 1 percentage point more likely to receive top grades, a gap that becomes statistically insignificant once school-level controls are included.

I find that although grade increases were often concentrated among socioeconomically disadvantaged students, gains in university placements were limited. The relationship between the degree of grade inflation and increases in the probability of placing into students’ top-choice programs follows an inverted-U: schools that most effectively converted grade gains into top-choice placements were concentrated in the middle of the school-quality distribution. Moreover, different grading channels worked unevenly: reassignment of top grades within a school was more strongly associated with placement gains at higher-quality schools, whereas broad, school-level grade inflation improved placement success mainly at lower-ranked schools.

My results qualitatively confirm previous evidence on institutional differences in grading patterns, providing the institutional context in which the well-documented student-level grading biases are relevant. Low quality schools took advantage of the grading discretion by employing a looser grading policy than high quality schools. A one unit lower quality difference between schools were associated with a 0.1 percent difference in the magnitude in grade inflation, with the higher quality school inflating less. This aligns with \cite{figlio2004}, \cite{10.1162/00335530360698441}, and \cite{dee2019} findings of stronger concentration of grade manipulation at lower quality schools. My results on gender biases quantitatively confirms the results of \cite{lavy2018} in which they find that female students are graded higher in non-blind exams. My results conflicts with \cite{hanna2012} and \cite{diamond} with neither paper finding evidence of gender bias. I claim the differences are driven by the differences in the grading procedure, where the two aforementioned papers examine grading bias on exams that students had already answered, while my results are on grading biases for predicted grades which potentially magnifies the grading biases as teachers are less informed about the students potential performance in exams.  Another explanation for the gender bias is the quality of the school. I find how gender bias was more concentrated in high quality schools and less in low quality schools. The results high light how institutional context shapes the differences in grading between gender groups.

I find how teachers grade students based on their economic affluence. I find that students from households with parents in higher income occupations received higher grades than students with parents in low-income occupations. The results align with \cite{hanna2012} in which they find how students from higher caste groups were graded higher than those at low caste groups. However, my results also find that the monotone relationship is largely observed in lower quality schools and less in high quality schools. Additionally I find that ethnic minorities were assigned lower grades than white students. This confirms the results from \cite{burgess2013}, in which they find how ethnic minorities are assigned lower grades in internal assessments.  

Although my results indicate a larger shift in grades for students in socio-economically disadvantaged backgrounds, the improvements in grades do not necessarily translate to an increase in an improvement in placements. This is due to socio-economically disadvantaged students often do not apply to top ranked universities albeit receiving a high grade in A-levels. On the other hand, students attending middle quality schools with good academic and socio-economic background leverage the upward shift in grades to secure placements at better universities. The inelasticity between grades and application to selective universities is similar to the results documented by \cite{NBERw18586}.  The upward shift of students from middle quality schools pulls students attending higher quality schools but with lower academic competence to lower their university placement. Despite positive relation between previous academic performance and the grade shift, the average academic score, measured by GCSE math, reduces at higher ranked universities. This indicates how non-academically related attributes, such as gender and socio-economic background dominates the upwards shifts due to academic performance. 

The paper is organized as follows: Section 2 provides the institutional background of the higher education system in the UK as well as the details of the policy set in place during the pandemic. Section 3 describes the data source I use in the analysis. Section 4 presents the empirical model and discuss the validity of the estimates along with the institutional details. Section 5 presents the results on heterogeneity of grading patterns across schools, as well as student level patterns, the dynamics of grading patterns, as well as their influence in placing students to their top choices. Section 6 discuss the magnitudes of the results with the literature on grading biases. Section 7 concludes with policy discussion on how my finding can inform policy makers with thinking about test-optional policies. 